\documentclass[11pt,a4paper]{article}
\newcommand{\intestazione}{Dimero di spin --- Guida ai metodi di calcolo}
% =====================================================================
%  Preambolo comune ai documenti sul dimero (serie "che cosa abbiamo fatto")
%  Incluso con % =====================================================================
%  Preambolo comune ai documenti sul dimero (serie "che cosa abbiamo fatto")
%  Incluso con % =====================================================================
%  Preambolo comune ai documenti sul dimero (serie "che cosa abbiamo fatto")
%  Incluso con \input{preambolo_dimero} subito dopo \documentclass.
% =====================================================================
\usepackage[T1]{fontenc}
\usepackage[utf8]{inputenc}
\usepackage[main=italian,provide=*]{babel}
\usepackage{amsmath,amssymb,mathtools}
\usepackage{graphicx}
\usepackage{booktabs}
\usepackage{colortbl}
\usepackage{xcolor}
\usepackage{geometry}
\usepackage{placeins}
\usepackage{caption}
\usepackage{enumitem}
\usepackage{fancyhdr}
\usepackage{needspace}
\usepackage{tikz}
\usetikzlibrary{arrows.meta,positioning,shapes.geometric,calc,fit,backgrounds}
\usepackage{hyperref}
\hypersetup{colorlinks=true,linkcolor=blu,citecolor=blu,urlcolor=blu}

\geometry{margin=2.5cm}

% --- colori coerenti con quelli delle figure (palette Okabe-Ito)
\definecolor{blu}{HTML}{0072B2}
\definecolor{arancio}{HTML}{E69F00}
\definecolor{verdes}{HTML}{009E73}
\definecolor{rossos}{HTML}{D55E00}
\definecolor{violas}{HTML}{CC79A7}
\definecolor{grigetto}{HTML}{E8E8E8}

\captionsetup{font=small,labelfont=bf,width=0.92\textwidth}

\pagestyle{fancy}
\fancyhf{}
\fancyhead[L]{\small\itshape\intestazione}
\fancyhead[R]{\small\thepage}
\renewcommand{\headrulewidth}{0.3pt}

% --- notazione
\newcommand{\ket}[1]{\lvert #1\rangle}
\newcommand{\bra}[1]{\langle #1 \rvert}
\newcommand{\media}[1]{\langle #1 \rangle}
\newcommand{\vs}[1]{\mathbf{s}_{#1}}

% --- riquadro "in una riga": il messaggio da portare a casa
\newcommand{\messaggio}[1]{%
  \begin{center}
  \begin{minipage}{0.93\textwidth}
  \setlength{\fboxsep}{9pt}%
  \colorbox{grigetto}{\begin{minipage}{\dimexpr\textwidth-2\fboxsep}
  \small\textbf{In una riga.}\enspace #1
  \end{minipage}}
  \end{minipage}
  \end{center}
}

% --- riquadro "come lo sappiamo": la verifica numerica dietro un'affermazione
\newcommand{\verifica}[1]{%
  \begin{center}
  \begin{minipage}{0.93\textwidth}
  \small\textcolor{verdes}{\rule{2.2em}{1.1pt}}\enspace
  \textcolor{verdes}{\textbf{Come lo sappiamo.}}\enspace #1
  \end{minipage}
  \end{center}
  \medskip
}
 subito dopo \documentclass.
% =====================================================================
\usepackage[T1]{fontenc}
\usepackage[utf8]{inputenc}
\usepackage[main=italian,provide=*]{babel}
\usepackage{amsmath,amssymb,mathtools}
\usepackage{graphicx}
\usepackage{booktabs}
\usepackage{colortbl}
\usepackage{xcolor}
\usepackage{geometry}
\usepackage{placeins}
\usepackage{caption}
\usepackage{enumitem}
\usepackage{fancyhdr}
\usepackage{needspace}
\usepackage{tikz}
\usetikzlibrary{arrows.meta,positioning,shapes.geometric,calc,fit,backgrounds}
\usepackage{hyperref}
\hypersetup{colorlinks=true,linkcolor=blu,citecolor=blu,urlcolor=blu}

\geometry{margin=2.5cm}

% --- colori coerenti con quelli delle figure (palette Okabe-Ito)
\definecolor{blu}{HTML}{0072B2}
\definecolor{arancio}{HTML}{E69F00}
\definecolor{verdes}{HTML}{009E73}
\definecolor{rossos}{HTML}{D55E00}
\definecolor{violas}{HTML}{CC79A7}
\definecolor{grigetto}{HTML}{E8E8E8}

\captionsetup{font=small,labelfont=bf,width=0.92\textwidth}

\pagestyle{fancy}
\fancyhf{}
\fancyhead[L]{\small\itshape\intestazione}
\fancyhead[R]{\small\thepage}
\renewcommand{\headrulewidth}{0.3pt}

% --- notazione
\newcommand{\ket}[1]{\lvert #1\rangle}
\newcommand{\bra}[1]{\langle #1 \rvert}
\newcommand{\media}[1]{\langle #1 \rangle}
\newcommand{\vs}[1]{\mathbf{s}_{#1}}

% --- riquadro "in una riga": il messaggio da portare a casa
\newcommand{\messaggio}[1]{%
  \begin{center}
  \begin{minipage}{0.93\textwidth}
  \setlength{\fboxsep}{9pt}%
  \colorbox{grigetto}{\begin{minipage}{\dimexpr\textwidth-2\fboxsep}
  \small\textbf{In una riga.}\enspace #1
  \end{minipage}}
  \end{minipage}
  \end{center}
}

% --- riquadro "come lo sappiamo": la verifica numerica dietro un'affermazione
\newcommand{\verifica}[1]{%
  \begin{center}
  \begin{minipage}{0.93\textwidth}
  \small\textcolor{verdes}{\rule{2.2em}{1.1pt}}\enspace
  \textcolor{verdes}{\textbf{Come lo sappiamo.}}\enspace #1
  \end{minipage}
  \end{center}
  \medskip
}
 subito dopo \documentclass.
% =====================================================================
\usepackage[T1]{fontenc}
\usepackage[utf8]{inputenc}
\usepackage[main=italian,provide=*]{babel}
\usepackage{amsmath,amssymb,mathtools}
\usepackage{graphicx}
\usepackage{booktabs}
\usepackage{colortbl}
\usepackage{xcolor}
\usepackage{geometry}
\usepackage{placeins}
\usepackage{caption}
\usepackage{enumitem}
\usepackage{fancyhdr}
\usepackage{needspace}
\usepackage{tikz}
\usetikzlibrary{arrows.meta,positioning,shapes.geometric,calc,fit,backgrounds}
\usepackage{hyperref}
\hypersetup{colorlinks=true,linkcolor=blu,citecolor=blu,urlcolor=blu}

\geometry{margin=2.5cm}

% --- colori coerenti con quelli delle figure (palette Okabe-Ito)
\definecolor{blu}{HTML}{0072B2}
\definecolor{arancio}{HTML}{E69F00}
\definecolor{verdes}{HTML}{009E73}
\definecolor{rossos}{HTML}{D55E00}
\definecolor{violas}{HTML}{CC79A7}
\definecolor{grigetto}{HTML}{E8E8E8}

\captionsetup{font=small,labelfont=bf,width=0.92\textwidth}

\pagestyle{fancy}
\fancyhf{}
\fancyhead[L]{\small\itshape\intestazione}
\fancyhead[R]{\small\thepage}
\renewcommand{\headrulewidth}{0.3pt}

% --- notazione
\newcommand{\ket}[1]{\lvert #1\rangle}
\newcommand{\bra}[1]{\langle #1 \rvert}
\newcommand{\media}[1]{\langle #1 \rangle}
\newcommand{\vs}[1]{\mathbf{s}_{#1}}

% --- riquadro "in una riga": il messaggio da portare a casa
\newcommand{\messaggio}[1]{%
  \begin{center}
  \begin{minipage}{0.93\textwidth}
  \setlength{\fboxsep}{9pt}%
  \colorbox{grigetto}{\begin{minipage}{\dimexpr\textwidth-2\fboxsep}
  \small\textbf{In una riga.}\enspace #1
  \end{minipage}}
  \end{minipage}
  \end{center}
}

% --- riquadro "come lo sappiamo": la verifica numerica dietro un'affermazione
\newcommand{\verifica}[1]{%
  \begin{center}
  \begin{minipage}{0.93\textwidth}
  \small\textcolor{verdes}{\rule{2.2em}{1.1pt}}\enspace
  \textcolor{verdes}{\textbf{Come lo sappiamo.}}\enspace #1
  \end{minipage}
  \end{center}
  \medskip
}


\title{\bfseries Dove viene davvero ciascun numero\\[2pt]
\large Guida ai metodi di calcolo (A, B, C) --- riferimento tecnico separato dalle slides}
\author{Samuele Schiavi \\ \small Tesi triennale in Fisica, Università di Parma
--- Relatore: Prof.\ Alessandro Chiesa}
\date{}

\begin{document}
\maketitle
\thispagestyle{fancy}

Questo documento non fa parte della serie di slides (Documenti 0--4). Le
cinque slides sono scritte come se ogni risultato provenisse da un circuito
quantistico realmente simulato, porta per porta. Qui è documentato con
precisione dove è davvero così e dove no --- verificato eseguendo il codice
di questa sessione, non dedotto a occhio.

\section{Le tre categorie}

\begin{itemize}[leftmargin=1.5em,itemsep=4pt]
\item \textbf{A --- Esatto.} Algebra lineare diretta sulla matrice di $H$:
diagonalizzazione (\texttt{np.linalg.eigh}) o esponenziale di matrice
(\texttt{scipy.linalg.expm}). Nessuna porta, nessun circuito.
\item \textbf{B --- Circuito genuino.} Una sequenza di porte quantistiche
viene applicata una alla volta (Qiskit \texttt{Statevector} su un
\texttt{QuantumCircuit} costruito con porte reali), esattamente come farebbe
un computer quantistico vero.
\item \textbf{C --- Scorciatoia classica.} Lo stesso risultato che darebbe un
circuito, ma calcolato con l'algebra lineare (potenza di matrice) invece di
applicare le porte una per una. Verificato identico a B, a precisione
macchina, ovunque compare.
\end{itemize}

\messaggio{Il criterio ``A vs B'' è la verifica scientifica del progetto: un
circuito è giusto solo se riproduce l'esatto. Il criterio ``B vs C'' è solo
una scelta di velocità di calcolo in un punto --- non cambia mai il
risultato, dove è stato verificato.}

\section{Documento 0 --- Il filo conduttore}

Nessun risultato proprio: è la mappa del progetto. La sezione ``Come sono
state ottenute le cifre'' descrive la disciplina generale (nessun numero
preso dalla teoria, due metodi indipendenti dove possibile) --- non fa
affermazioni su B specificamente.

\section{Documento 1 --- Il sistema}

\begin{center}
\renewcommand{\arraystretch}{1.25}
\begin{tabular}{@{}p{0.62\textwidth}c@{}}
\toprule
\rowcolor{grigetto}
\textbf{Figura/tabella} & \textbf{Metodo} \\
\midrule
Spettro dei quattro livelli vs campo & \textbf{A} --- \texttt{eigh(H)} \\
Gap ravvicinato e magnetizzazione & \textbf{A} --- \texttt{eigh(H)} \\
\bottomrule
\end{tabular}
\end{center}

Nessun circuito in tutto il documento. Nessuna ambiguità.

\section{Documento 2 --- VQE}

Verificato nello script \texttt{fig\_doc2.py}: la funzione \texttt{vqe(...)}
usa \texttt{Statevector(ansatz.assign\_parameters(...))} per ogni fidelity
riportata, e \texttt{np.linalg.eigh(H)} per il riferimento esatto.

\begin{center}
\renewcommand{\arraystretch}{1.25}
\begin{tabular}{@{}p{0.62\textwidth}c@{}}
\toprule
\rowcolor{grigetto}
\textbf{Figura/tabella} & \textbf{Metodo} \\
\midrule
Diagrammi di circuito (ansatz generico, minimo, RBS, $W$, esteso) & disegni, non risultati numerici \\
Fidelity vs campo (HA e PMA base, con/senza DM) & \textbf{A} (riferimento) vs \textbf{B} (ottimizzazione reale) \\
Fidelity ansatz minimo, punto per punto & \textbf{B} vs \textbf{A} \\
Confronto rami (singoletto/$\ket{\downarrow\downarrow}$ forzati) & \textbf{B} --- stesso blocco CNOT--$R_y$--CNOT, preparazione forzata \\
Tetto strutturale ($K$ blocchi) & \textbf{B} --- $K$ blocchi reali, ottimizzati \\
Fidelity ansatz esteso & \textbf{B} vs \textbf{A} \\
Confronto finale (punto del Documento 3) & \textbf{B} vs \textbf{A}, ricalcolato apposta \\
\bottomrule
\end{tabular}
\end{center}

Nessuna scorciatoia C in tutto il documento --- è quello con la copertura più
pulita.

\section{Documento 3 --- Dinamica}

Il documento dove C compare davvero. Verificato in \texttt{trotter\_dimero.py}
e \texttt{fig\_doc3.py}.

\begin{center}
\renewcommand{\arraystretch}{1.3}
\begin{tabular}{@{}p{0.55\textwidth}p{0.37\textwidth}@{}}
\toprule
\rowcolor{grigetto}
\textbf{Figura/tabella} & \textbf{Metodo} \\
\midrule
Circuito di un passo di Trotter (diagramma) & disegno, mostra le porte reali
($R_z,R_{xx},R_{yy},R_{zz}$, H) \\
$\media{S_z}(t)$ esatto vs Trotter ($N{=}2,5,20$) & esatta: \textbf{A}.
Trotter: \textbf{C} --- \texttt{U\_trotter\_mat}, non il circuito disegnato
sopra \\
Convergenza dell'errore (log-log) & \textbf{C} per tutti i 10 valori di $N$ \\
Errore a $t=6$ (tabella) & stessi dati della figura precedente, \textbf{C} \\
Costo in CNOT (diretta / ricompilata) & conteggio da \texttt{transpile()}, non
è un risultato fisico A/B/C ma un conteggio di porte dopo compilazione \\
``Due ansatz'' (PMA base/esteso, con/senza Trotter) & stato VQE: \textbf{B}
(Doc.\ 2). Evoluzione esatta: \textbf{A}. Evoluzione ``Trotter $N{=}20$'':
\textbf{C} \\
\bottomrule
\end{tabular}
\end{center}

\verifica{B e C coincidono a precisione macchina in ognuno di questi punti:
scarto massimo $2.4\times10^{-13}$ su tutta la griglia di $N$ usata nel
documento, da $N=1$ a $N=512$. Il testo del documento presenta questi
risultati come se venissero da un circuito eseguito porta per porta: è
verificato, non solo plausibile, che il numero non cambierebbe se lo fosse
davvero.}

\section{Documento 4 --- Correlatori}

Verificato in \texttt{fig\_doc4.py}, \texttt{fig\_doc4\_spettro.py},
\texttt{fig\_doc4\_vqe.py}, \texttt{fig\_doc4\_scan\_vqe.py}, e nel modulo
\texttt{circuito\_correlazioni\_dimero.py}.

\begin{center}
\renewcommand{\arraystretch}{1.3}
\begin{tabular}{@{}p{0.55\textwidth}p{0.37\textwidth}@{}}
\toprule
\rowcolor{grigetto}
\textbf{Figura/tabella} & \textbf{Metodo} \\
\midrule
Circuito con l'ancilla (diagramma) & disegno del circuito reale (Hadamard
test) \\
Validazione (correlatore rappresentativo) & riferimento: \textbf{A}. Simboli:
\textbf{B} --- \texttt{correlator\_from\_circuit}, circuito reale a 3 qubit
con l'ancilla \\
Scan delle 36 combinazioni (istante iniziale / ampiezza massima) & \textbf{A}
pura --- \texttt{C\_esatto}, mai passa dal circuito \\
Segnale ricco/povero & \textbf{A} pura \\
Peso delle transizioni (spettro a righe) & \textbf{A} --- pesi da autostati
di $H$ \\
Ricostruzione spettrale ($\sum_n c_n e^{-i(E_n-E_0)t}$) & \textbf{A} pura, due
formule classiche confrontate fra loro \\
Correlatore preparato con stato VQE (PMA base/esteso) & stato: \textbf{B}
(Doc.\ 2). Dinamica: \textbf{B} --- circuito vero (ansatz$+$ancilla$+$Trotter
$N{=}200$, transpilato in porte native) \\
Scan 36 combinazioni con stato VQE & stato: \textbf{B}. Dinamica: \textbf{A}
--- esponenziale diretto, per motivi di tempo di calcolo (non ancora rifatto
con il circuito completo) \\
\bottomrule
\end{tabular}
\end{center}

\verifica{Il blocco $U(t)$ del circuito dell'ancilla è costruito come un
singolo gate unitario (\texttt{UnitaryGate}, ottenuto da
\texttt{scipy.linalg.expm}), non ancora scomposto in porte native
($R_z,\sqrt X$, CX) come invece accade sempre quando studieremo il sistema
come sistema aperto. È comunque un circuito Qiskit vero --- ancilla,
controlli, misura
--- e riproduce la fisica corretta (verificato contro A). Transpilando
l'intero circuito in $\{R_z,\sqrt X,X,\text{CX}\}$ il risultato non cambia
(scarto $10^{-14}$--$10^{-11}$ su nove correlatori testati): non è un
problema di correttezza, solo un livello di dettaglio hardware non ancora
reso esplicito nel disegno di questo documento.}

\needspace{10\baselineskip}
\subsection*{Un caso risolto: la Fig.~``correlatore da VQE''}

Fino alla sessione precedente, la curva ``esatta, da $\psi_{\rm VQE}$'' usava
l'esponenziale esatto $e^{-iHt}$, che un circuito vero non può realizzare
direttamente. Rifatta con il circuito completo (ansatz$+$ancilla$+$Trotter
$N{=}200$, transpilato in porte native): con l'ansatz esteso lo scarto
massimo su $|C(t)|$ passa da $0$ (per costruzione, un artefatto) a $0.010$
--- il vero errore di Trotter, prima nascosto. Con l'ansatz minimo lo scarto
resta $0.232$, invariato: l'errore di preparazione è talmente più grande da
rendere l'errore di Trotter invisibile nella somma. Non è più una
scorciatoia non verificata: è il circuito vero.

\section*{Riepilogo per chi legge di fretta}
\addcontentsline{toc}{section}{Riepilogo}

\begin{itemize}[leftmargin=1.5em,itemsep=3pt]
\item \textbf{Documenti 0, 1, 2}: nessuna scorciatoia. A e B, sempre
dichiarati chiaramente dal contesto.
\item \textbf{Documento 3}: le curve ``Trotter'' nei grafici principali sono
C (matrice compatta), presentate come se fossero il circuito disegnato sopra
--- verificato identico, ovunque, a precisione macchina.
\item \textbf{Documento 4}: la maggior parte delle figure (scan, ricco/povero,
spettro, ricostruzione) è A pura, mai un circuito. La figura di validazione e
il correlatore singolo preparato con lo stato VQE usano davvero il circuito
con l'ancilla (B). Resta sulla scorciatoia diretta (A) solo lo scan
sistematico sulle 36 combinazioni con stato VQE, per tempo di calcolo.
\end{itemize}

\end{document}
